\section{Critical Gaps and Limitations of Previous Studies}

Prior work leaves three connected gaps for event-sensitive urban air-quality
forecasting. First, temporal forecasters and numerical retrieval systems do not
necessarily preserve the provenance, coverage, and availability assumptions of
local contextual records. A numerically similar PM2.5 window can occur under a
different weather or event regime, and a retrospective event label can be
mistaken for information available at forecast time. The SACB addresses this
gap by joining official records under explicit source, coverage, and
availability audits.

Second, historical retrieval can introduce look-ahead bias when a candidate
input or target overlaps the query context. Similarity alone does not prevent a
near-duplicate future from entering the retrieved evidence. The LSER addresses
this gap through a strict temporal embargo and training-history restriction,
followed by separately auditable PM2.5, weather, calendar, and event similarity
terms.

Third, retrieval gains are often reported without testing whether a strong local
context model already captures the useful signal, whether the gain persists
under distribution shift, or whether an explanation can be traced to stored
evidence. The DFEH addresses this gap by evaluating direct XGBoost and frozen
Chronos-Bolt paths on common origins, exposing shift indicators, and retaining
machine-readable feature, retrieval, event, and uncertainty evidence. This
component supports comparison and auditability; it does not convert contextual
association into causal explanation.

These mappings delimit the novelty of the study. Event-TimeRAF does not contain
a learned dual-encoder retriever or internal Channel Prompting, and the present
single-county experiment does not establish geographic generalization.

\section{Core Contributions}

The proposed framework contributes three mechanisms corresponding one-to-one
with the gaps above:
\begin{itemize}
    \item \textit{Source-Audited Context Builder (SACB):} constructs an aligned
    85-dimensional context from official PM2.5, meteorological, calendar, and
    event records while retaining source and availability metadata.
    \item \textit{Leakage-Safe Event-Context Retriever (LSER):} retrieves
    training-history analogues only after enforcing a candidate-target embargo,
    then exposes the contribution of each similarity channel and the retrieved
    trajectory evidence.
    \item \textit{Drift-Aware Forecast and Evidence Head (DFEH):} places direct
    horizon models, frozen-TSFM forecast fusion, drift-state analysis, and
    traceable explanation records within one chronological evaluation contract.
\end{itemize}

The empirical contribution is intentionally qualified: weather-calendar
context produces the clearest supported gain, whereas event and retrieval
components provide selective rather than uniformly positive improvements.
